\chapter{Harvard Professor Explains The Rules of Writing}

This first lecture in the course is not mathematical in the literal blackboard sense, but it is rigorous. Pinker gives us a causal model of bad prose, a practical method for revision, and a sequence of increasingly fine distinctions: between expertise and communication, between abstraction and imagery, between speech and writing, between clarity and pleasure, and finally between human prose and the smooth generic competence of large language models. What follows is a set of companion notes that keeps the lecture's order and treats its central claims as a system we can actually use.

\section{Bad Writing and the Curse of Knowledge}

The lecture opens with a useful reversal. The easy explanation of bad writing is moral: perhaps academics, bureaucrats, and professionals choose obscurity because obscurity signals status. Pinker does not deny that such motives exist, but he insists that they do not explain the general case. Many bad writers are intelligent, sincere, and full of substantive ideas. The failure is more basic: they cannot reliably recover the state of mind of the person who does not already know what they know.

\begin{definition}
The curse of knowledge is the difficulty of reconstructing what it is like not to know something that we ourselves know.
\end{definition}

We can summarize Pinker's mechanism schematically as
\begin{equation}
\text{expert knowledge} \to \text{failed audience model} \to \text{jargon / acronym / abstraction} \to \text{reader confusion}.
\end{equation}
The lecture earns this formula through examples rather than through definition alone. The molecular biologist at a TED conference loses the room in a few seconds because he speaks as if everyone shares the internal setting of his field. Acronyms go unexplained. The problem being solved is never stated. The talk begins in the middle.

The same mistake shows up in prose that says something like ``the level of the stimulus was proportional to the intensity of the reaction'' when the real content is much simpler: children look longer at a bunny than at a truck. The sentence is not false. It is opaque.

\subsection{Question \& Answer}

\paragraph{Question.} Why is there so much bad writing if the writers are often intelligent and sincere?

\paragraph{Answer.} Because intelligence in a field is not the same as skill at locating the reader. Pinker's point is that the writer usually fails not by having nothing to say, but by forgetting what the audience cannot infer. The first enemy is not malice. It is miscalibration.

\begin{proposition}
If we had to name a single dominant flaw in explanatory prose, Pinker would locate it in the writer's failure to model the reader's mind.
\end{proposition}

This is why he links the curse of knowledge to egocentrism and to the absence of a working theory of mind. The lecture begins, then, not with style in the narrow sense, but with a cognitive mistake that infects every later choice of wording.

\section{Audience Calibration as a Writing Practice}

Once the cause has been identified, the lecture turns immediately to method. If the failure is a bad audience model, then the remedy cannot be only introspection. Pinker says that he does try to imagine the reader, but he also says, very plainly, that imagination is not enough. One does not know, from the inside, when one is under the spell of the curse of knowledge.

The revision loop is
\begin{align}
\text{draft} &\to \text{imagined reader}, \\
\text{draft} + \text{actual reader response} &\to \text{revision}.
\end{align}
We first do what we can by empathy. We then test that empathy against actual readers.

The lecture is especially good here because it refuses a lazy picture of the ``ordinary reader.'' Pinker did not give drafts to his mother because she was unsophisticated. On the contrary, he stresses that she was intelligent and well read. What mattered was that she was not a cognitive psychologist. The relevant reader is not ``the average person.'' The relevant reader is an intellectually serious person who does not already inhabit the writer's local specialization.

We should preserve the sequence of readers that the lecture supplies:
\begin{itemize}
\item a trusted intelligent non-specialist,
\item a professional editor,
\item academics from other fields,
\item even near neighbors inside the university who are still outside the tiny circle that generated the draft.
\end{itemize}
The striking claim is that even adjacent academics often fail to understand one another because each lab, subfield, or method-group builds up a private atmosphere of shorthand.

This gives us a practical rule. Before we refine a sentence for elegance, we should ask who has actually read it and where that reader stands relative to our field. The chapter's later advice about examples and rhythm rests on this earlier discipline.

\section{Concreteness, Imagery, and Sensory Understanding}

The next pivot is crucial. The lecture does not remain at the level of audience-testing. Pinker widens the frame and asks what understanding itself consists in. His answer is that language is a means, not the endpoint. Readers do not understand by merely receiving one verbal string after another. They understand by building some grasp of an idea, and that grasp is very often visual, motoric, emotional, auditory, or otherwise sensory.

Schematically,
\begin{equation}
\text{language} \to \text{mental model},
\qquad
\text{not merely}
\qquad
\text{language} \to \text{more language}.
\end{equation}

That is why abstraction is dangerous when it arrives too early. A term like ``stimulus'' may be efficient within a research paper, but it gives the general reader almost nothing to inspect. A bunny and a truck do.

\subsection{Worked Example}

Pinker's own contrast gives us a clean miniature derivation from abstraction to observability:
\begin{enumerate}
\item Start with the sentence ``the level of the stimulus was proportional to the intensity of the reaction.''
\item Notice that the sentence names a relation but hides the scene.
\item Ask what the entities actually are.
\item Replace the compressed terms by the concrete case: children, bunny, truck, looking longer.
\item Recover a statement the reader can picture and therefore understand.
\end{enumerate}

We can write the contrast schematically as
\begin{align}
\text{``The level of the stimulus was proportional to the intensity of the reaction.''}
&\not\Rightarrow \text{a stable mental scene},\\
\text{``Kids look longer at a bunny than at a truck.''}
&\Rightarrow \text{a picture the reader can inspect}.
\end{align}

From here the lecture moves naturally to visual metaphor. Earlier writers often sound more vivid, Pinker suggests, because they had to reach for shared images before academic abstraction had become so readily available. ``Aggression'' is efficient. ``The spirit of the hawk kneaded into our flesh'' is memorable because it lets the reader see something. We need not imitate older prose mechanically, but we should understand the function it was performing.

\section{Why Writing Is Harder Than Speaking}

At this point the interviewer introduces a new tension: if speech comes so naturally, why does writing remain slow and effortful? Pinker's answer is not romantic. He does not say that writing is harder because it is more noble. He says it is harder because speech is supported by structures that writing lacks.

\subsection{Question \& Answer}

\paragraph{Question.} Why does writing feel unnatural compared with speech?

\paragraph{Answer.} Because conversation begins inside a shared situation, whereas writing must manufacture its own context on the page.

The lecture's contrast can be stated compactly as
\begin{align}
\text{conversation} &= \text{common ground} + \text{deixis} + \text{feedback},\\
\text{writing} &= \text{text on the page}.
\end{align}

This short formula gathers several claims that Pinker develops one by one:
\begin{enumerate}
\item In conversation, the speakers already know why they are there.
\item They can use words like ``this,'' ``that,'' ``she,'' and ``the thing'' because the setting supplies the missing structure.
\item They receive correction in real time through faces, questions, pauses, and visible drift of attention.
\item In writing, the reader may be distant in place, distant in time, and wholly outside the social setting in which the sentences were produced.
\item Therefore the page must carry burdens that speech can offload onto context.
\end{enumerate}

This beat is essential to the lecture's rhythm. Without it, the advice about examples would sound like a stylistic preference. With it, the advice becomes a necessity. Writing must replace lost common ground by explicit structure.

\section{Generalization, Example, and the Shape of Meaning}

The lecture now becomes methodological in a narrower sense. Pinker endorses the claim that generalizations without examples and examples without generalizations both fail, though he softens the word ``useless.'' The point is not that either form is worthless in every setting. The point is that by themselves they do not complete the act of understanding.

We may state the balance as
\begin{equation}
\text{examples} + \text{generalizations} \to \text{understanding}.
\end{equation}
The interviewer reformulates this as a balance between context and compression, and that summary is useful as long as we remember it is a schematic gloss:
\begin{equation}
\text{context} \leftrightarrow \text{compression}.
\end{equation}

The logic is worth making explicit.
\begin{enumerate}
\item A generalization suppresses detail; it sweeps over particulars.
\item Without a case, the reader may not know which particular region of experience is being named.
\item An example pins down the ballpark.
\item But an example alone leaves the reader asking, ``And your point is?''
\item A generalization then states what the example exemplifies.
\end{enumerate}

The lecture's own examples are especially revealing because they move from explanatory prose into word meaning itself. Bathroom does not literally require a bath. Breakfast need not be a literal breaking of a fast. Christmas need not be understood compositionally as Christ's Mass each time the word is used. Olive oil and baby oil illustrate different semantic relations inside apparently similar compounds. The Lederer examples go further by showing that even familiar morphology is historically layered: singer is transparently built by a live rule, whereas finger is not.

\begin{remark}
Meaning is often conventional, historical, and only partly compositional. A reader cannot be expected to recover an intended sense from bare form alone when the language itself does not always behave that way.
\end{remark}

This is why the lecture lingers over etymology and fossilized rules. The point is not antiquarian. The point is that writers should not assume that the surface structure of a phrase automatically delivers understanding.

\section{Style as Pleasure: Rhythm, Sound, and Brevity}

After clarity comes pleasure. The lecture is careful about this order. Style is not introduced as ornament applied after meaning is settled. Rather, once we have made prose intelligible, we still have the question of whether it moves well, sounds well, and rewards the reader.

\subsection{Question \& Answer}

\paragraph{Question.} What makes prose not just clear but pleasurable?

\paragraph{Answer.} Pinker's answer is that good prose has audible structure. It can be spoken. It carries rhythm. It avoids needless friction. It may use alliteration or other sound-patterns lightly, not as decoration pasted on afterward, but as part of the sentence's ease and force.

The recommended local test is simple:
\begin{equation}
\text{read aloud} \to \text{hear friction} \to \text{revise cadence}.
\end{equation}
If we cannot say a sentence smoothly, the reader will probably not move through it smoothly either.

Several variables enter here:
\begin{itemize}
\item \textbf{metrical structure}: prose has beats even when it is not verse;
\item \textbf{sibilance}: too many successive hissing sounds can roughen the surface;
\item \textbf{alliteration}: a little can help a sentence glide;
\item \textbf{brevity}: compression removes both cognitive load and stylistic drag.
\end{itemize}

Pinker treats ``brevity is the soul of wit'' and Strunk's ``omit needless words'' as more than slogans. They are examples of themselves. There is no excess in the way the rule is stated. The lecture also gives a practical argument for brevity: when a fixed word limit forces a piece to contract, the prose often improves rather than merely shrinks.

We may summarize that claim as
\begin{equation}
\text{needless words} \to \text{extra processing} \to \text{weaker force}.
\end{equation}

Humor belongs here for the same reason. A joke that arrives late loses force. A sentence that signals its punch line too heavily steps on itself. The lecture's account of wit is therefore continuous with its account of prose style more generally: the shorter line often lands harder, provided it remains apt.

\section{Old Prose, Academic Prose, and LLM Prose}

The final movement of the lecture gathers the earlier pieces into a wider historical frame. Why does older prose so often sound vivid? Pinker gives several reasons in sequence. Many writers were trained on strong literary models. Prose served more openly as a medium of self-presentation. Most important for the logic of the lecture, earlier writers had fewer ready-made abstractions and fewer dead cliches available to them, so they were pushed back toward image, metaphor, and concrete locution.

This historical explanation leads directly into Pinker's attack on bad academic prose. The harm is not only aesthetic. It is wasted thought, wasted effort, and preventable misunderstanding. When a grant review, student paper, or manuscript must be read five or six times to recover the claim, the prose has failed in an ethical as well as stylistic sense: it has consumed attention without delivering clarity.

By the end of the lecture, this entire framework is tested against LLM writing. Pinker grants its strengths before criticizing it. That order matters. LLM prose is often plain, orderly, and structurally competent. It tends to produce introductions, conclusions, and sensible sequencing. But it also tends toward the generic and the banal.

Schematically,
\begin{align}
\text{LLM prose} &= \text{clear ordering} + \text{plain syntax} + \text{generic pastiche},\\
\text{strong prose} &= \text{clarity} + \text{concreteness} + \text{freshness}.
\end{align}

So we should not confuse clarity with liveliness:
\begin{equation}
\text{clear prose} \neq \text{fresh prose}.
\end{equation}
Pinker closes by widening the intellectual implication once more. If large models can extract so much from enormous bodies of text, then any theory of mind that gives pride of place only to explicit rules and logical programming is incomplete. Yet he stops well short of identifying human intelligence with the large language model. Human learning is situated, interactive, embodied, and developmentally timed in a way that the model is not.

\section{Summary}

The lecture gives us a disciplined sequence rather than a bag of tips. We begin with the curse of knowledge, because bad prose is first a failure to locate the reader. We then move to actual revision practice, because empathy must be tested. From there we learn why writing needs concreteness: understanding is not a stream of words but a constructed mental model. We then see why writing is harder than speaking, why examples and generalizations must cooperate, why meaning is often less compositional than we imagine, and why style matters as rhythm, sound, brevity, and pleasure. The closing discussion of older prose, academic prose, and LLM prose does not replace the earlier lessons. It reframes them. The central task remains the same: we must say what we mean in a form that another mind can actually enter.